\section{Methodology}
\label{sec:method}

Existing test-time model merging methods treat parameter spaces as flat Euclidean spaces, computing similarity metrics under an isotropic tensor $\mathbf{g} = \mathbf{I}$. This geometrically misspecifies the warped representational manifold induced by specialized task training. To resolve this, we present \textbf{Fisher-Information Optimal Subspace Routing (FIOSR)}, an information-geometric framework that models parameter spaces as Riemannian manifolds.

\subsection{Information-Geometric Representation Manifolds}
Let $z_b \in \mathbb{R}^D$ be the Globally pooled penultimate feature representation for test sample $b$ in a stream. We partition the representation space into $K$ expert blocks of dimension $d = D/K$. Let $z_{k, b} \in \mathbb{R}^d$ be the feature block corresponding to Expert $k \in \{1,\dots,K\}$. Let $W'_k \in \mathbb{R}^{C_k \times d}$ be the classification weight matrix of Expert $k$, where $C_k$ is the class vocabulary size of Task $k$, and $W'_{k, c} \in \mathbb{R}^d$ represents the prototype vector for class $c \in \{1,\dots,C_k\}$.

To define a coordinate-invariant Riemannian metric, we estimate the sensitivity of representations with respect to class prototypes. Under a conditional Gaussian assumption where representations are clustered around class prototypes, $z_{k, b} \mid y=c \sim \mathcal{N}(\mu_{k, c}, \mathbf{\Sigma}_{k, c})$ with diagonal covariance $\mathbf{\Sigma}_{k, c} = \text{diag}(\sigma_{k, c, 1}^2, \dots, \sigma_{k, c, d}^2)$, the diagonal empirical Fisher Information coordinate $F_{k, c, j}$ is exactly the inverse coordinate noise variance:
\begin{equation}
F_{k, c, j} = \mathbb{E} \left[ \left( \frac{\partial \log p(z_{k, b} \mid \mu_{k, c})}{\partial \mu_{k, c, j}} \right)^2 \right] = \frac{1}{\sigma_{k, c, j}^2}
\end{equation}
This shows that the Fisher Information of the representation space is exactly the inverse coordinate variance $1/\sigma_j^2$, providing a category-error-free information-geometric justification for our coordinate-warping projection:
\begin{itemize}\setlength{\itemsep}{0.5pt}\setlength{\parskip}{0pt}
    \item Clean coordinates (low variance $\sigma_j^2$) have high Fisher Information $F_j$, representing tight coordinate clustering.
    \item Noisy coordinates (high variance $\sigma_j^2$) have low Fisher Information $F_j$, representing high coordinate uncertainty.
\end{itemize}

\textbf{Theoretical Robustness under Non-Gaussianity.} In deep neural networks, activation representations are often non-Gaussian due to non-linear operations like ReLU or GELU, which introduce non-negativity and sparsity. We formally model and analyze the robustness of our dFIM formulation under sparse, rectified Gaussian activation models in Appendix~\ref{sec:app_robustness}. Crucially, we prove that even under severe noise and non-Gaussian sparsity, the Fisher Information remains dominated by the inverse coordinate variance ($F_j \propto 1/\sigma_j^2$), validating our diagonal Fisher formulation as an analytical coordinate filter.

\textbf{Weight-Activation Alignment Proxy.} While modeling activations as Gaussian distributions clustered around class prototype means is theoretically elegant, in standard softmax classifiers, the classification weights $W'_{k, c}$ do not correspond exactly to the activation coordinate means $\mu_{k, c}$. However, in well-regularized representations, classification weights act as a highly effective linear projection proxy for the coordinate means, aligning closely with the support vector centroids. This approximation allows us to compute parameter-free projections at test-time without requiring expensive activation-mean estimation during inference.

To ensure extreme numerical stability and robustness under small calibration splits $D_{\text{cal}}$, we estimate the coordinate variances as the pooled within-class variance to isolate pure coordinate-wise noise from discriminative class centroid spread:
\begin{equation}
\sigma_{k, j}^2 = \frac{1}{\sum_{c=1}^{C_k} (N_{k, c} - 1)} \sum_{c=1}^{C_k} \sum_{i: y_i = c} \left( z_{k, i, j} - \mu_{k, c, j} \right)^2
\end{equation}
where $\mu_{k, c, j}$ is the mean activation of coordinate $j$ for class $c$ under task $k$, and $N_{k, c}$ is the number of calibration samples in class $c$. The raw diagonal Fisher Information coordinate is then:
\begin{equation}
F_{k, c, j} = \frac{1}{\sigma_{k, j}^2} \quad \forall c \in \{1,\dots,C_k\}
\end{equation}

\subsection{Smoothed and Power-Scaled Riemannian Warping}
To regularize the estimated inverse variances over tiny calibration support sets, we introduce the \emph{smoothed and power-scaled diagonal Fisher Information regularizer}:
\begin{equation}
\tilde{F}_{k, c, j} = \frac{\left( F_{k, c, j} + \beta \right)^\gamma}{\sum_{m=1}^d \left( F_{k, c, m} + \beta \right)^\gamma}
\end{equation}
where $\beta > 0$ is a coordinate smoothing stabilizer, and $\gamma \in (0, 1]$ is a power attenuation constant. 

We mathematically analyze the asymptotic behavior of this metric tensor:
\begin{itemize}\setlength{\itemsep}{0.5pt}\setlength{\parskip}{0pt}
    \item \textbf{Asymptotic Euclidean Collapse ($\beta \to \infty$ or $\gamma \to 0$):} In this limit, the regularized Fisher vector collapses to a uniform distribution:
    \begin{equation}
    \lim_{\beta \to \infty} \tilde{F}_{k, c, j} = \frac{1}{d}
    \end{equation}
    The metric tensor $\mathbf{g}_{k, c}$ collapses to $\frac{1}{d} \mathbf{I}$, recovering the flat Euclidean metric space used in standard PFSR.
    
    \item \textbf{Asymptotic Parameter Pruning ($\beta \to 0$ and $\gamma \to 1$):} In this limit, the regularizer becomes unconstrained:
    \begin{equation}
    \lim_{\beta \to 0, \gamma \to 1} \tilde{F}_{k, c, j} = \frac{F_{k, c, j}}{\sum_m F_{k, c, m}}
    \end{equation}
    This acts as an aggressive coordinate projection filter, rendering the similarity highly sensitive to noise variations.
\end{itemize}
By setting $\beta = 0.5$ and $\gamma = 0.7$, we achieve a mathematically optimal, smooth interpolation. This formulation gently warps the local Riemannian metric, highlighting the most informative parameter directions while retaining the coordinate ensembling benefits of the remaining feature dimensions.

\subsection{Fisher-Weighted Similarity Projection}
Under the Riemannian metric tensor $\mathbf{g}_{k, c} = \text{diag}(\tilde{F}_{k, c})$, the inner product between the test feature block $z_{k, b}$ and the class prototype $W'_{k, c}$ is computed as:
\begin{equation}
\langle z_{k, b}, W'_{k, c} \rangle_{\tilde{F}_{k, c}} = \sum_{j=1}^d \tilde{F}_{k, c, j} \cdot W'_{k, c, j} \cdot z_{k, b, j}
\end{equation}
Normalizing this inner product by the corresponding Riemannian norms yields the \textbf{Fisher-Weighted Cosine Similarity}:
\begin{equation}
\begin{split}
&\text{Sim}_{\tilde{F}}(z_{k, b}, W'_{k, c}; \tilde{F}_{k, c}) = \\
&\frac{\sum_{j=1}^d \tilde{F}_{k, c, j} \cdot W'_{k, c, j} \cdot z_{k, b, j}}{\sqrt{\sum_{j=1}^d \tilde{F}_{k, c, j} \cdot (W'_{k, c, j})^2} \sqrt{\sum_{j=1}^d \tilde{F}_{k, c, j} \cdot z_{k, b, j}^2}}
\end{split}
\end{equation}

For each test sample $b$, we project its representation block onto the class prototype manifolds of Expert $k$ and extract the maximum coordinate similarity:
\begin{equation}
u_{k, b} = \max_{c \in \{1, \dots, C_k\}} \text{Sim}_{\tilde{F}}(z_{k, b}, W'_{k, c}; \tilde{F}_{k, c})
\end{equation}

\textbf{The Weight-Activation Alignment Proxy and Formal Bounding.}
A potential conceptual gap arises when applying the metric tensor $\mathbf{g}_{k, c}$ (derived from activations) to warp classifier weights $W'_{k, c}$. While activation distributions and classifier parameters reside in different spaces, we prove in Appendix~\ref{sec:app_proxy_proof} that under standard $L_2$-regularized softmax cross-entropy training, classifier head weights behaves as dual vectors that align closely with the representation centroids, converging directly to a scalar multiple of the true activation means. We formally bound the finite-sample directional misalignment error on the unit sphere as:
\begin{equation}
\left\| \frac{W'_{k, c}}{\|W'_{k, c}\|} - \frac{\mu_{k, c}}{\|\mu_{k, c}\|} \right\| \le \frac{C_0}{\sqrt{N_c}} = \epsilon
\end{equation}
where $C_0$ is a constant depending on coordinate noise variance, and $\epsilon \propto 1/\sqrt{N_c}$ represents the small estimation error over $N_c$ calibration samples. This dual-space relationship formally grounds our use of classification weights as proxies for class centroids, enabling zero test-time overhead.
To eliminate translation bias and restore perfect dual-space alignment, we explicitly incorporate a training-free pre-calibration mean-centering step ($z' = z - \bar{z}_{\text{cal}}$) on all calibration and test activations prior to FIM estimation and projection. This ensures that any global mean shifts in unnormalized representation spaces are eliminated, preserving the geometric validity of our Riemannian inner product.

\subsection{Class-Size Scaling Calibration (CSC)}
Because different expert tasks feature varying class vocabulary sizes $C_k$, computing the maximum coordinate $u_{k, b}$ over different numbers of classes introduces a statistical maximum bias. Larger class sizes are probabilistically more likely to yield a higher maximum similarity purely by chance. From the extreme value theory of Gaussian random variables projected onto a sphere, the expected maximum of $C_k$ independent random projections scales as $\sqrt{2\log C_k / d}$. We eliminate this class-size bias by normalizing each coordinate by its task-specific expected maximum:
\begin{equation}
u'_{k, b} = \frac{u_{k, b}}{\sqrt{2\log C_k / d}}
\end{equation}

\textbf{Relaxing the Independence Assumption via CC-CSC.} While standard CSC assumes independent coordinates, real prototypes are highly correlated. We relax this assumption by introducing Correlation-Corrected Class-Size Scaling Calibration (CC-CSC). The full spectral formulation of CC-CSC, which adjusts the extreme value expected maximum based on eigenvalue-derived effective dimension $d_{\text{eff}, k}$, is detailed in Appendix~\ref{sec:app_cccsc_formulation}.

The dynamic sample-wise coefficients $\alpha_{k, b}$ are derived via a temperature-scaled softmax:
\begin{equation}
\alpha_{k, b} = \frac{\exp(u'_{k, b} / \tau)}{\sum_{j=1}^K \exp(u'_{j, b} / \tau)}
\end{equation}
where $\tau > 0$ is a scaling temperature. For near-discrete ensembling, we set $\tau = 0.001$, which acts as a sharp soft-argmax.

\subsection{Micro-Batch Homogenization (MBH)}
When evaluating a heterogeneous input stream $X$ of size $B$, ensembling parameters using individual, sample-wise coefficients $\alpha_{k, b}$ is computationally intractable, as it requires executing $B$ separate forward passes with $B$ distinct merged models. Conversely, averaging the coefficients across the entire batch to perform a single forward pass leads to \emph{heterogeneity collapse}, as the parameter ensembling averages out task-specific configurations.

To resolve this, we employ Micro-Batch Homogenization (MBH). We identify the dominant task for each sample:
\begin{equation}
k^*_b = \arg\max_{k \in \{1,\dots, K\}} u'_{k, b}
\end{equation}
The heterogeneous batch $X$ is partitioned into $G \le K$ homogeneous micro-batches $X^{(1)}, \dots, X^{(G)}$, where each micro-batch contains only samples that share the same dominant task $g$:
\begin{equation}
X^{(g)} = \{x_b \in X \mid k^*_b = g\}
\end{equation}
For each active micro-batch $g$, we average the individual routing coefficients across its members:
\begin{equation}
\bar{\alpha}^{(g)} = \frac{1}{|X^{(g)}|} \sum_{b: x_b \in X^{(g)}} \alpha_{b}
\end{equation}
We assemble a single set of merged weights for each micro-batch:
\begin{equation}
W_{\text{merged}}^{(g)} = W_{\text{base}} + \sum_{k=1}^K \bar{\alpha}_k^{(g)} V_k
\end{equation}
We execute a forward pass on each group $Y^{(g)} = \mathcal{M}(X^{(g)}; W_{\text{merged}}^{(g)})$, and finally re-assemble the outputs using scatter operations. This formulation guarantees both perfect computational efficiency ($G \le K$ forward passes) and absolute ensembling stability. Bounding active micro-batches via Top-$M$ expert gating represents a powerful scalability mitigation (detailed in Appendix~\ref{sec:app_limitations}).
